\documentclass[11pt]{article}
\usepackage[margin=1in]{geometry}
\usepackage{amsmath,amssymb,bm}
\usepackage[T1]{fontenc}
\usepackage[utf8]{inputenc}
\usepackage{lmodern}
\usepackage{microtype}
\usepackage{booktabs}
\usepackage{graphicx}
\usepackage{enumitem}
\setlist[itemize]{leftmargin=1.5em}

\title{SAVLG-CBM Branch Architectures}
\author{}
\date{\today}

\begin{document}
\maketitle

\section{Purpose}
This note describes the two branch architectures used in the current
\texttt{SAVLG-CBM} implementation:
\begin{itemize}
    \item \textbf{shared branch architecture}
    \item \textbf{dual branch architecture}
\end{itemize}
The key difference is whether the model uses one shared concept head for both
global concept prediction and spatial grounding, or whether it uses separate
global and spatial concept heads.

\section{Common SAVLG Pipeline}
For an image $x$, the visual backbone produces spatial feature tensors. Let
\[
\mathbf{f}_5 \in \mathbb{R}^{C_5 \times H_5 \times W_5}
\]
denote the final-stage feature map, and optionally let
\[
\mathbf{f}_4 \in \mathbb{R}^{C_4 \times H_4 \times W_4}
\]
denote an earlier feature map when a multiscale branch is used.

For $K$ concepts, SAVLG produces:
\begin{itemize}
    \item a global concept logit vector
    \[
    \mathbf{c}_{\mathrm{global}} \in \mathbb{R}^{K},
    \]
    \item a spatial concept map tensor
    \[
    \mathbf{M} \in \mathbb{R}^{K \times H \times W},
    \]
    \item a pooled spatial concept vector
    \[
    \mathbf{c}_{\mathrm{spatial}} \in \mathbb{R}^{K},
    \]
    \item and the fused concept logits
    \[
    \mathbf{c}_{\mathrm{final}}
    =
    \mathbf{c}_{\mathrm{global}}
    +
    \alpha \, \mathbf{c}_{\mathrm{spatial}}.
    \]
\end{itemize}

In the current implementation, the image-level concept BCE is applied to
$\mathbf{c}_{\mathrm{final}}$, so the spatial branch participates directly in
concept supervision through the residual fusion term.

\section{Architecture 1: Shared Branch}
\subsection{Definition}
In the shared architecture, a single concept layer produces the spatial concept
maps:
\[
\mathbf{M} = s(\mathbf{f}),
\qquad
\mathbf{M} \in \mathbb{R}^{K \times H \times W}.
\]
There is no separate global concept head. Instead, the same output tensor is
used for both roles:
\begin{itemize}
    \item the spatial path keeps $\mathbf{M}$ as the localization map
    \item the global path pools the same tensor to produce concept logits
\end{itemize}

Concretely, if the concept layer does not expose a separate
\texttt{forward\_both(...)} method, the implementation reuses the same tensor
for both outputs:
\[
\mathbf{G} = \mathbf{M},
\qquad
\mathbf{S} = \mathbf{M},
\]
where $\mathbf{G}$ denotes the tensor consumed by the global pooling path and
$\mathbf{S}$ denotes the tensor consumed by the spatial path.
Then the global concept logits are obtained by pooling the same map tensor:
\[
\mathbf{c}_{\mathrm{global}}
=
\operatorname{Pool}_{\mathrm{global}}(\mathbf{M}).
\]

\subsection{Computation Graph}
The shared architecture therefore has the form:
\[
x \rightarrow \text{backbone} \rightarrow \mathbf{f}
\rightarrow s(\cdot) \rightarrow \mathbf{M}
\]
\[
\mathbf{c}_{\mathrm{global}}
=
\operatorname{Pool}_{\mathrm{global}}(\mathbf{M}),
\qquad
\mathbf{c}_{\mathrm{spatial}}
=
\operatorname{Pool}_{\mathrm{spatial}}(\mathbf{M}).
\]
Finally,
\[
\mathbf{c}_{\mathrm{final}}
=
\mathbf{c}_{\mathrm{global}}
+
\alpha \mathbf{c}_{\mathrm{spatial}}.
\]

\subsection{Interpretation}
This architecture forces one concept tensor to satisfy two jobs at once:
\begin{itemize}
    \item \textbf{semantic discrimination}: the pooled map must predict concept
    presence well
    \item \textbf{spatial faithfulness}: the map itself should align with the
    annotated concept region
\end{itemize}

The advantage is simplicity. The disadvantage is coupling: the same parameters
must optimize global prediction and localization simultaneously, which can make
the representation less specialized.

\section{Architecture 2: Dual Branch}
\subsection{Definition}
In the dual architecture, the model has two distinct heads:
\begin{itemize}
    \item a \textbf{global head} for concept prediction
    \item a \textbf{spatial head} for concept localization
\end{itemize}

The implementation exposes this explicitly through:
\[
\texttt{forward\_global(...)} \quad \text{and} \quad \texttt{forward\_spatial(...)}.
\]
The forward pass returns
\[
\texttt{global\_outputs}, \texttt{spatial\_maps}
=
\texttt{forward\_both(...)}.
\]

\subsection{Split-Stage Dual Variant}
The basic dual variant keeps the global path on conv5 features and allows the
spatial path to use a configurable stage:
\[
\mathbf{c}_{\mathrm{global}} = g(\mathbf{f}_5),
\qquad
\mathbf{M} = s(\mathbf{f}_{\mathrm{spatial}}),
\]
where
\[
\mathbf{f}_{\mathrm{spatial}} \in \{\mathbf{f}_5, \mathbf{f}_4, \ldots\}.
\]
Thus, the global and spatial branches are separate even if they may read from
the same or nearby backbone stages.

\subsection{Multiscale conv4+conv5 Dual Variant}
The most important dual specialization in this codebase is the
\texttt{multiscale\_conv45} branch.

\paragraph{Global path.}
The global branch stays on conv5:
\[
\mathbf{c}_{\mathrm{global}} = g(\mathbf{f}_5).
\]
When the VLG-style global head is used, this is effectively:
\[
\bar{\mathbf{f}}_5 = \operatorname{GAP}(\mathbf{f}_5),
\qquad
\mathbf{c}_{\mathrm{global}} = W_g \bar{\mathbf{f}}_5 + b_g.
\]

\paragraph{Spatial path.}
The spatial branch first fuses conv4 and conv5 features:
\[
\tilde{\mathbf{f}}_4 = W_4 * \mathbf{f}_4,
\qquad
\tilde{\mathbf{f}}_5 = \operatorname{Up}(W_5 * \mathbf{f}_5),
\]
\[
\mathbf{f}_{\mathrm{fused}}
=
\operatorname{ReLU}(\tilde{\mathbf{f}}_4 + \tilde{\mathbf{f}}_5),
\]
and then applies the spatial concept head:
\[
\mathbf{M} = s(\mathbf{f}_{\mathrm{fused}}).
\]

So the dual multiscale model uses:
\begin{itemize}
    \item a dedicated \textbf{global conv5 head}
    \item a dedicated \textbf{spatial multiscale conv4+conv5 head}
\end{itemize}

\subsection{Computation Graph}
The dual architecture has the form:
\[
x \rightarrow \text{backbone} \rightarrow (\mathbf{f}_4,\mathbf{f}_5)
\]
\[
\mathbf{c}_{\mathrm{global}} = g(\mathbf{f}_5),
\qquad
\mathbf{M} = s(\mathbf{f}_{\mathrm{spatial}} \text{ or } \mathbf{f}_{\mathrm{fused}}),
\]
\[
\mathbf{c}_{\mathrm{spatial}}
=
\operatorname{Pool}_{\mathrm{spatial}}(\mathbf{M}),
\qquad
\mathbf{c}_{\mathrm{final}}
=
\mathbf{c}_{\mathrm{global}} + \alpha \mathbf{c}_{\mathrm{spatial}}.
\]

\subsection{Interpretation}
The dual architecture allows specialization:
\begin{itemize}
    \item the global branch can remain close to a strong image-level concept
    predictor, often warm-started from VLG
    \item the spatial branch can focus on localization quality
    \item the two branches are combined only at the concept-logit level
\end{itemize}

This separation is especially important when:
\begin{itemize}
    \item the user wants to freeze the global head
    \item the global branch should preserve VLG semantics
    \item the spatial branch should use a different stage or a multiscale fusion
\end{itemize}

\section{Global and Spatial Pooling}
Once the spatial map tensor $\mathbf{M}$ is produced, the model pools it to
obtain spatial concept logits. In the residual path, the implementation uses
log-sum-exp pooling:
\[
c_{\mathrm{spatial},k}
=
\tau \log \sum_{u} \exp\!\left(M_{k,u}/\tau\right)
- \tau \log |U|.
\]
This acts as a smooth max over locations.

For the global path:
\begin{itemize}
    \item if the head already returns a vector, no extra pooling is needed
    \item if the head returns maps, those maps are pooled to produce
    $\mathbf{c}_{\mathrm{global}}$
\end{itemize}

\section{Main Conceptual Difference}
The core distinction is:
\begin{center}
\begin{tabular}{@{}p{0.24\linewidth}p{0.33\linewidth}p{0.33\linewidth}@{}}
\toprule
Property & Shared & Dual \\
\midrule
Global head & implicit, reused from spatial maps & explicit dedicated head \\
Spatial head & same tensor as global path & separate branch \\
Specialization & weak & strong \\
Supports VLG-style global warm-start cleanly & limited & yes \\
Supports multiscale spatial fusion with separate global path & awkward & yes \\
\bottomrule
\end{tabular}
\end{center}

\section{Key Differences in Practice}
The architectural difference is not just a code-level refactor. It changes what
the model is allowed to specialize.

\subsection{Shared Branch}
\begin{itemize}
    \item One concept tensor must support both image-level concept prediction
    and localization.
    \item The pooled global concept logits are obtained from the same tensor that
    is later interpreted as the spatial explanation map.
    \item This is simpler, but it couples semantic discrimination and spatial
    faithfulness tightly.
\end{itemize}

\subsection{Dual Branch}
\begin{itemize}
    \item The global branch has its own head and can stay close to a strong VLG
    concept predictor.
    \item The spatial branch has its own head and can use a different feature
    source or a multiscale fused feature map.
    \item This decouples classification quality from localization quality and
    makes freezing or warm-starting the global branch much cleaner.
\end{itemize}

\subsection{Why this mattered empirically}
In the recent \texttt{resnet18} reproduction work, the main recovery did not
come from the outside penalty alone. The bigger gain came from restoring the
older branch factorization:
\begin{itemize}
    \item \texttt{branch\_arch=dual}
    \item \texttt{spatial\_branch\_mode=multiscale\_conv45}
    \item VLG-style global warm-start
    \item optional freezing of the global head
\end{itemize}
That combination recovered the stronger sparse results much more clearly than
the newer shared-branch runs.

\section{Observed Performance}
Table~\ref{tab:branch-performance} summarizes the most relevant comparison from
the recent \texttt{resnet18} experiments.

\begin{table}[h]
\centering
\begin{tabular}{@{}p{0.34\linewidth}ccccc@{}}
\toprule
Run & Branch style & Outside & Full acc & ACC@5 & AVGACC \\
\midrule
Recent weak \texttt{conv5} run & shared & 0.0 & 0.7561 & 0.6929 & 0.7362 \\
Old-style matched run & dual + multiscale\_conv45 & 0.0 & 0.7591 & 0.7566 & 0.7574 \\
Old-style matched run & dual + multiscale\_conv45 & 0.25 & 0.7591 & 0.7577 & 0.7585 \\
\bottomrule
\end{tabular}
\caption{Representative \texttt{resnet18 conv5} comparison. The main gap is not
explained by outside penalty alone. Restoring the old dual-branch recipe
recovers most of the lost sparse performance.}
\label{tab:branch-performance}
\end{table}

\paragraph{Interpretation.}
The full accuracies are all close, but the sparse accuracies are not. The
shared-branch run stays around
\[
\text{AVGACC} \approx 0.736,
\]
whereas the old-style dual-branch reproductions recover
\[
\text{AVGACC} \approx 0.758.
\]
So the main practical difference is:
\begin{itemize}
    \item \textbf{shared} is simpler but underperforms on sparse concept
    prediction
    \item \textbf{dual} preserves a stronger global concept predictor while
    allowing a specialized spatial branch
\end{itemize}

\section{Why the Dual Architecture Performed Better}
The recent reproduction work suggests that the stronger older results were not
explained by outside penalty alone. The more important change was restoring the
older \emph{dual} branch recipe:
\begin{itemize}
    \item separate global and spatial heads
    \item multiscale conv4+conv5 spatial fusion
    \item VLG-style global warm-start
    \item optional freezing of the global head
\end{itemize}
In other words, the main gain came from restoring a better branch factorization,
while the outside penalty provided only a smaller additional improvement.

\section{Summary}
\begin{itemize}
    \item The \textbf{shared} architecture uses one concept map tensor for both
    global and spatial roles.
    \item The \textbf{dual} architecture explicitly separates the global concept
    path from the spatial localization path.
    \item The multiscale conv4+conv5 model is a specialized \textbf{dual}
    architecture: global on conv5, spatial on fused conv4+conv5.
    \item The dual design better supports warm-starting, freezing, and branch
    specialization, which matches the empirical recovery of the stronger older
    \texttt{resnet18} results.
\end{itemize}

\section{Backbone Result Tables}
This appendix collects the current headline performance and localization numbers
for the two backbones used most heavily in this project. The localization
protocols are not identical across all rows, so each table explicitly states the
metric family being used.

\subsection{ResNet-18: Cross-Method Summary}
\begin{table}[h]
\centering
\small
\resizebox{\textwidth}{!}{
\begin{tabular}{@{}lccccccc@{}}
\toprule
Method & Full acc & ACC@5 & AVGACC & mean IoU & mAP@0.3 & mAP@0.5 & mAP@0.7 \\
\midrule
LF-CBM & 0.7402 & 0.5147 & 0.6817 & 0.0750 & 0.0366 & 0.0361 & 0.0364 \\
SALF-CBM & 0.7320 & 0.5335 & 0.6826 & 0.0789 & 0.0306 & 0.0295 & 0.0287 \\
VLG-CBM & 0.7594 & 0.7546 & 0.7556 & 0.0874 & 0.1054 & 0.0491 & 0.0288 \\
SAVLG-CBM (best residual, $\alpha=0.20$) & 0.7606 & 0.7585 & 0.7585 & 0.1062 & 0.1174 & 0.0546 & 0.0313 \\
\bottomrule
\end{tabular}
}
\caption{ResNet-18 cross-method summary. Localization numbers here are the
matched post-hoc Grad-CAM comparison from \texttt{docs/cub\_results.md}.}
\label{tab:r18-cross-method}
\end{table}

\begin{table}[h]
\centering
\small
\resizebox{\textwidth}{!}{
\begin{tabular}{@{}lcccccc@{}}
\toprule
Method & mean IoU & mAP@0.3 & mAP@0.5 & mAP@0.7 & point hit & coverage \\
\midrule
LF-CBM (Grad-CAM) & 0.0750 & 0.0366 & 0.0361 & 0.0364 & 0.5626 & 0.0334 \\
SALF-CBM (Grad-CAM) & 0.0789 & 0.0306 & 0.0295 & 0.0287 & 0.7284 & 0.0446 \\
VLG-CBM (Grad-CAM) & 0.0874 & 0.1054 & 0.0491 & 0.0288 & 0.8000 & 0.2536 \\
SAVLG-CBM (Grad-CAM) & 0.1062 & 0.1174 & 0.0546 & 0.0313 & 0.9011 & 0.1341 \\
\midrule
SALF-CBM (native) & 0.0778 & 0.00645 & 0.00500 & 0.00456 & -- & -- \\
SAVLG-CBM (native) & 0.2820 & 0.4268 & 0.1790 & 0.0356 & -- & -- \\
\bottomrule
\end{tabular}
}
\caption{\texttt{Resnet18} localization comparison. The first four rows are the
strict matched \texttt{gt\_present} mean-threshold Grad-CAM comparison across
\texttt{LF}, \texttt{SALF}, \texttt{VLG}, and \texttt{SAVLG}. The final two
rows append the best available native-map references for \texttt{SALF} and
\texttt{SAVLG}, so they should be read as native reference rows rather than as
part of the matched Grad-CAM protocol.}
\label{tab:r18-gradcam-localization}
\end{table}

\begin{table}[h]
\centering
\small
\resizebox{\textwidth}{!}{
\begin{tabular}{@{}lccccccl@{}}
\toprule
Method / run & mean IoU & mAP@0.3 & mAP@0.5 & mAP@0.7 & point hit & coverage & Notes \\
\midrule
SALF-CBM (native pairwise artifact) & 0.0778 & 0.00645 & 0.00500 & 0.00456 & -- & -- & native SALF vs SAVLG comparator artifact \\
SAVLG-CBM (native clean-global pairwise artifact) & 0.0233 & 0.00620 & 0.00505 & 0.00458 & -- & -- & paired with the SALF native artifact above \\
SAVLG-CBM (best residual \texttt{conv4+conv5}) & 0.1309 & 0.1189 & 0.0656 & 0.0311 & 0.9889 & 0.3657 & strongest earlier native \texttt{resnet18} row from the residual line \\
SAVLG-CBM (exact matched native rerun, \texttt{conv4+conv5}) & 0.2820 & 0.4268 & 0.1790 & 0.0356 & -- & -- & old native evaluator family, reproduced exactly \\
\bottomrule
\end{tabular}
}
\caption{\texttt{Resnet18} native-map localization references. \texttt{SALF}
and \texttt{SAVLG} do not have a single fully matched native-map table covering
all later \texttt{SAVLG} variants, so we separate the earlier paired
\texttt{SALF}/clean-global-\texttt{SAVLG} artifact from the stronger later
\texttt{SAVLG} residual native localizers. This table is useful for native-map
context, but it is not a single strict four-method benchmark in the way
Table~\ref{tab:r18-gradcam-localization} is.}
\label{tab:r18-native-localization}
\end{table}

\subsection{ResNet-18: SAVLG Internal Ablations}
\begin{table}[h]
\centering
\small
\resizebox{\textwidth}{!}{
\begin{tabular}{@{}lccccccc@{}}
\toprule
SAVLG variant & Full acc & ACC@5 & AVGACC & mean IoU & mAP@0.3 & mAP@0.5 & mAP@0.7 \\
\midrule
Soft-align + outside (warm-start frozen-global) & 0.7591 & 0.7549 & 0.7560 & 0.1246 & 0.0491 & 0.0336 & 0.0237 \\
Residual-Alpha-0.05 & 0.7592 & 0.7560 & 0.7568 & 0.1309 & 0.1189 & 0.0656 & 0.0311 \\
Residual-Alpha-0.10 & 0.7585 & 0.7572 & 0.7581 & 0.1307 & 0.1189 & 0.0657 & 0.0310 \\
Residual-Alpha-0.05 + local-MIL & 0.7591 & 0.7563 & 0.7575 & 0.1216 & 0.0790 & 0.0434 & 0.0250 \\
Residual conv5-only & 0.7592 & 0.7560 & 0.7567 & 0.1251 & 0.1168 & 0.0630 & 0.0312 \\
Residual conv4-only & 0.7587 & 0.7547 & 0.7559 & 0.1445 & 0.1170 & 0.0710 & 0.0395 \\
Old-style dual conv5, outside=0.25 & 0.7591 & 0.7577 & 0.7585 & pending & pending & pending & pending \\
Old-style dual conv5, outside=0 & 0.7591 & 0.7566 & 0.7574 & 0.0543 & 0.0281 & 0.0034 & 0.00047 \\
\bottomrule
\end{tabular}
}
\caption{ResNet-18 SAVLG internal ablations. Native localization numbers are
reported where they are already available; the old-style outside=0.25
full-test localization rerun was interrupted before completion.}
\label{tab:r18-savlg}
\end{table}

\subsection{ResNet-50: Backbone Summary}
\begin{table}[h]
\centering
\small
\resizebox{\textwidth}{!}{
\begin{tabular}{@{}lcccccc@{}}
\toprule
Method / run & Full acc & ACC@5 & AVGACC & Loc. metric @0.3 & Loc. metric @0.5 & Notes \\
\midrule
VLG-CBM (resnet50\_cub\_mm) & 0.8529 & 0.8501 & 0.8533 & 0.1373 & 0.0528 & Grad-CAM sweep \\
SALF-CBM (resnet50\_cub\_mm) & 0.7263 & 0.2936 & 0.5540 & 0.0066 & 0.0051 & Grad-CAM mean-threshold \\
SAVLG conv5 + cov025 & 0.8566 & 0.8558 & 0.8564 & 0.3301 & 0.1408 & MaxBoxAcc sweep \\
SAVLG conv5 + cov050 & 0.8566 & 0.8544 & 0.8561 & 0.3207 & 0.1371 & MaxBoxAcc sweep \\
SAVLG conv5 + dice010 & 0.8570 & 0.8558 & 0.8563 & 0.3393 & 0.1462 & MaxBoxAcc sweep \\
SAVLG conv4 & 0.8584 & 0.8551 & 0.8559 & 0.3709 & 0.1573 & MaxBoxAcc sweep \\
SAVLG conv4+conv5 & 0.8573 & 0.8565 & 0.8561 & 0.4467 & 0.2152 & best localization \\
SAVLG conv5 soft-align only & 0.8615 & 0.8570 & 0.8568 & pending & pending & strongest conv5 classifier-side result \\
\bottomrule
\end{tabular}
}
\caption{ResNet-50 summary. SAVLG rows use native-map sweep metrics. VLG and
SALF rows use the best available ResNet-50 Grad-CAM localization artifacts from
the pod results directory, so this table mixes protocols and should be read as
a broad comparison rather than a strictly matched localization benchmark.}
\label{tab:r50-summary}
\end{table}

\begin{table}[h]
\centering
\small
\resizebox{\textwidth}{!}{
\begin{tabular}{@{}lccccl@{}}
\toprule
Run & mean IoU & mAP@0.3 & mAP@0.5 & Full acc & Notes \\
\midrule
VLG-CBM (Grad-CAM, resnet50\_cub\_mm) & 0.1091 & 0.1373 & 0.0528 & 0.8529 & best available pod-side Grad-CAM sweep \\
SALF-CBM (Grad-CAM, resnet50\_cub\_mm) & 0.0789 & 0.0066 & 0.0051 & 0.7263 & best available pod-side Grad-CAM artifact \\
SAVLG conv5 + cov025 & 0.1824 & 0.3301 & 0.1408 & 0.8566 & baseline conv5 branch \\
SAVLG conv5 + cov050 & 0.1793 & 0.3207 & 0.1371 & 0.8566 & stronger coverage did not help \\
SAVLG conv5 + dice010 & 0.1854 & 0.3393 & 0.1462 & 0.8570 & best finished conv5 localization \\
SAVLG conv4 & 0.2437 & 0.3709 & 0.1573 & 0.8584 & best dense run in the 5-run sweep \\
SAVLG conv4+conv5 & 0.2822 & 0.4467 & 0.2152 & 0.8573 & best localization in the branch sweep \\
SAVLG conv5 soft-align only & 0.1241 & pending & pending & 0.8615 & strongest completed conv5 classifier-side run \\
SAVLG soft-align + outside penalty & 0.1246 & -- & -- & 0.7591 & frozen-global warm-start local-loss ablation \\
SAVLG soft-align + local-MIL + outside & 0.1255 & -- & -- & 0.7591 & best reported ResNet-50 mean IoU \\
\bottomrule
\end{tabular}
}
\caption{Dedicated ResNet-50 SAVLG localization table. For the branch/loss sweep
rows, mean IoU is the best value achieved over the explicit activation-threshold
sweep, and the `mAP@0.3` / `mAP@0.5` columns correspond to the logged
MaxBoxAcc sweep values at those IoU thresholds. The added VLG/SALF rows come
from the best available pod-side Grad-CAM artifacts, so this table is useful
for broad comparison but is not a perfectly matched single-protocol benchmark.}
\label{tab:r50-localization-diagnostics}
\end{table}

\begin{table}[h]
\centering
\small
\resizebox{0.88\textwidth}{!}{
\begin{tabular}{@{}lccccc@{}}
\toprule
Backbone & Checkpoint & mean IoU & mAP@0.3 & mAP@0.5 & Full acc \\
\midrule
\texttt{resnet18\_cub} & SAVLG dual \texttt{conv4+conv5} & 0.2820 & 0.4268 & 0.1790 & 0.7592 \\
\texttt{resnet50\_cub\_mm} & SAVLG dual \texttt{conv4+conv5} & 0.2822 & 0.4469 & 0.2152 & 0.8573 \\
\bottomrule
\end{tabular}
}
\caption{Exact matched native-map comparison for the best SAVLG dual \texttt{conv4+conv5}
checkpoint on \texttt{resnet18} and \texttt{resnet50}. Both rows use the same
old native evaluator family, \texttt{map\_source=native},
\texttt{map\_normalization=concept\_zscore\_minmax}, fixed threshold sweep
\texttt{0.05..0.95}, \texttt{concept\_mode=intersection}, and
\texttt{eval\_subset\_mode=gt\_present}. This is the clean apples-to-apples
localization comparison between the two backbones.}
\label{tab:conv45-r18-vs-r50-exact}
\end{table}

\subsection{Compact Cross-Backbone Method Comparison}
\begin{table}[h]
\centering
\small
\resizebox{\textwidth}{!}{
\begin{tabular}{@{}llccccc@{}}
\toprule
Backbone & Method & Full acc & ACC@5 & AVGACC & Loc.@0.3 & Loc.@0.5 \\
\midrule
\texttt{resnet18\_cub} & LF-CBM & 0.7402 & 0.5147 & 0.6817 & 0.0366 & 0.0361 \\
\texttt{resnet18\_cub} & SALF-CBM & 0.7320 & 0.5335 & 0.6826 & 0.0306 & 0.0295 \\
\texttt{resnet18\_cub} & VLG-CBM & 0.7594 & 0.7546 & 0.7556 & 0.1054 & 0.0491 \\
\texttt{resnet18\_cub} & SAVLG-CBM (best residual) & 0.7606 & 0.7585 & 0.7585 & 0.1174 & 0.0546 \\
\midrule
\texttt{resnet50\_cub\_mm} & LF-CBM & N/A & N/A & N/A & N/A & N/A \\
\texttt{resnet50\_cub\_mm} & SALF-CBM & 0.7263 & 0.2936 & 0.5540 & 0.0066 & 0.0051 \\
\texttt{resnet50\_cub\_mm} & VLG-CBM & 0.8529 & 0.8501 & 0.8533 & 0.1373 & 0.0528 \\
\texttt{resnet50\_cub\_mm} & SAVLG-CBM (best localization) & 0.8573 & 0.8565 & 0.8561 & 0.4467 & 0.2152 \\
\bottomrule
\end{tabular}
}
\caption{Compact cross-backbone comparison using only the headline method rows.
For \texttt{resnet18}, localization columns are the matched Grad-CAM comparison.
For \texttt{resnet50}, SAVLG localization uses native-map MaxBoxAcc sweep
results, while VLG and SALF use the best available pod-side Grad-CAM
artifacts. This table intentionally omits auxiliary SAVLG ablations such as
local-MIL and intermediate residual variants.}
\label{tab:compact-cross-backbone}
\end{table}

\subsection{All Methods Across Both Backbones}
\begin{table}[h]
\centering
\small
\resizebox{\textwidth}{!}{
\begin{tabular}{@{}llccccc@{}}
\toprule
Backbone & Method & Full acc & ACC@5 & AVGACC & Loc.@0.3 & Loc.@0.5 \\
\midrule
\texttt{resnet18\_cub} & LF-CBM & 0.7402 & 0.5147 & 0.6817 & 0.0366 & 0.0361 \\
\texttt{resnet18\_cub} & SALF-CBM & 0.7320 & 0.5335 & 0.6826 & 0.0306 & 0.0295 \\
\texttt{resnet18\_cub} & VLG-CBM & 0.7594 & 0.7546 & 0.7556 & 0.1054 & 0.0491 \\
\texttt{resnet18\_cub} & SAVLG-CBM & 0.7606 & 0.7585 & 0.7585 & 0.1174 & 0.0546 \\
\midrule
\texttt{resnet50\_cub\_mm} & LF-CBM & N/A & N/A & N/A & N/A & N/A \\
\texttt{resnet50\_cub\_mm} & SALF-CBM & 0.7263 & 0.2936 & 0.5540 & 0.0066 & 0.0051 \\
\texttt{resnet50\_cub\_mm} & VLG-CBM & 0.8529 & 0.8501 & 0.8533 & 0.1373 & 0.0528 \\
\texttt{resnet50\_cub\_mm} & SAVLG-CBM & 0.8573 & 0.8565 & 0.8561 & 0.4467 & 0.2152 \\
\bottomrule
\end{tabular}
}
\caption{All currently tracked headline method rows across the two main
backbones. For \texttt{resnet18}, localization values are the matched
cross-method Grad-CAM comparison. For \texttt{resnet50}, SAVLG localization
values are native-map MaxBoxAcc sweep results, while VLG and SALF values come
from the best available pod-side Grad-CAM artifacts.}
\label{tab:all-methods-both-backbones}
\end{table}

\subsection{Residual Pooling Ablation}
\begin{table}[h]
\centering
\small
\resizebox{\textwidth}{!}{
\begin{tabular}{@{}lcccccc@{}}
\toprule
Pooling & Full acc & ACC@5 & AVGACC & Best mean IoU & Best Loc.@0.3 & Best Loc.@0.5 \\
\midrule
\texttt{lse} & 0.7591 & 0.7577 & 0.7585 & pending & pending & pending \\
\texttt{avg} & 0.7606 & 0.7561 & 0.7576 & 0.05435 & 0.02805 & 0.00331 \\
\texttt{topk (0.2)} & 0.7613 & 0.7566 & 0.7579 & 0.05435 & 0.02805 & 0.00331 \\
\bottomrule
\end{tabular}
}
\caption{Residual pooling ablation for the old-style dual \texttt{conv5}
\texttt{resnet18} recipe with outside penalty. \texttt{avg} and
\texttt{topk(0.2)} native-map localization were evaluated with an explicit
threshold sweep.}
\label{tab:pooling-ablation}
\end{table}

\paragraph{Main takeaways.}
\begin{itemize}
    \item On \texttt{resnet18}, the main step change came from restoring the
    older dual-branch residual recipe, not from outside penalty alone.
    \item On \texttt{resnet18}, residual \texttt{conv4}-only gives the strongest
    native localization among the spatial-source ablations, while
    \texttt{conv5}-only and multiscale remain strongest on sparse accuracy.
    \item On \texttt{resnet50}, \texttt{conv4+conv5} is the clear localization
    winner, while \texttt{conv5 + soft-align only} is the strongest completed
    conv5 classifier-side run.
\end{itemize}

\end{document}
